
\documentclass{article}
\usepackage{colortbl}
\usepackage{makecell}
\usepackage{multirow}
\usepackage{supertabular}

\begin{document}

\newcounter{utterance}

\centering \large Interaction Transcript for game `referencegame', experiment `random\_grids', episode 5 with reflexion{-}llama{-}3.1{-}8b{-}instruct.
\vspace{24pt}

{ \footnotesize  \setcounter{utterance}{1}
\setlength{\tabcolsep}{0pt}
\begin{supertabular}{c@{$\;$}|p{.15\linewidth}@{}p{.15\linewidth}p{.15\linewidth}p{.15\linewidth}p{.15\linewidth}p{.15\linewidth}}
    \# & $\;$A & \multicolumn{4}{c}{Game Master} & $\;\:$B\\
    \hline

    \theutterance \stepcounter{utterance}  
    & & \multicolumn{4}{p{0.6\linewidth}}{
        \cellcolor[rgb]{0.9,0.9,0.9}{
            \makecell[{{p{\linewidth}}}]{
                \texttt{\tiny{[P1$\langle$GM]}}
                \texttt{You have attempted this same task before and did not succeed. These are the plans you wrote after those attempts.} \\
\texttt{Trial \#1: New plan: I failed to distinguish the target grid from the distractors because my description was too general and didn't take into account the unique features of the target grid. The target grid has a row with three X's in the first, third, and fifth columns, but this description also applies to the first row of the first distractor grid. Next time, I will focus on a more specific feature that is unique to the target grid, such as the position of the X's in the second row or the arrangement of the □'s in the third row.} \\
\\ 
\texttt{Trial \#2: New plan: I failed to distinguish the target grid from the distractors because my description was too general and didn't take into account the unique features of the target grid. The target grid has a row with three X's in the first, third, and fifth columns, but this description also applies to the first row of the first distractor grid. However, I didn't consider the arrangement of the □'s in the third row, which is unique to the target grid. Next time, I will focus on a more specific feature that is unique to the target grid, such as the arrangement of the □'s in the third row, and combine it with the row with three X's in the second row. For example, I could describe the target grid as "The grid that has a row with three X's in the first, third, and fifth columns of the second row, and a row with two □'s in the first and second columns of the third row."} \\
\\ 
\texttt{Trial \#3: New plan: I failed to distinguish the target grid from the distractors because my description was too general and didn't take into account the unique features of the target grid. The target grid has a row with three X's in the first, third, and fifth columns, but this description also applies to the first row of the first distractor grid. However, I didn't consider the arrangement of the □'s in the third row, which is unique to the target grid. Next time, I will focus on a more specific feature that is unique to the target grid, such as the arrangement of the □'s in the third row, and combine it with the row with three X's in the second row. For example, I could describe the target grid as "The grid that has a row with three X's in the first, third, and fifth columns of the second row, and a row with two □'s in the first and second columns of the third row."} \\
\\ 
\\ 
\texttt{You are given three grids, where each of them is 5 by 5 in size.} \\
\texttt{Grids have empty cells marked with "▢" and filled cells marked with "X".} \\
\texttt{Your task is to generate a referring expression that best describes the target grid while distinguishing it from the two other distractor grids.} \\
\texttt{The first grid is the target grid, and the following two grids are the distractors.} \\
\\ 
\texttt{Target grid:} \\
\\ 
\texttt{X X □ □ X} \\
\texttt{□ X X □ X} \\
\texttt{X □ X X □} \\
\texttt{□ X □ X X} \\
\texttt{X □ □ X X} \\
\\ 
\texttt{Distractor grid 1:} \\
\\ 
\texttt{X X □ X □} \\
\texttt{□ X X □ X} \\
\texttt{X □ □ X □} \\
\texttt{□ X □ □ X} \\
\texttt{X □ X □ X} \\
\\ 
\texttt{Distractor grid 2:} \\
\\ 
\texttt{□ X □ X □} \\
\texttt{X □ X □ X} \\
\texttt{□ X □ X □} \\
\texttt{X □ X □ X} \\
\texttt{□ X □ X □} \\
\\ 
\texttt{Instruction: Describe the target grid.} \\
\texttt{Generate the referring expression starting with the tag "Expression: " for the given target grid. Omit any other text.} \\
            }
        }
    }
    & \\ \\

    \theutterance \stepcounter{utterance}  
    & \multicolumn{4}{p{0.6\linewidth}}{
        \cellcolor[rgb]{0.9,0.9,0.9}{
            \makecell[{{p{\linewidth}}}]{
                \texttt{\tiny{[P1$\rangle$GM]}}
                \texttt{Expression: The grid that has a row with three X's in the first, third, and fifth columns of the second row, and a row with two X's in the first and third columns of the fourth row.} \\
            }
        }
    }
    & & \\ \\

    \theutterance \stepcounter{utterance}  
    & & & \multicolumn{2}{p{0.3\linewidth}}{
        \cellcolor[rgb]{0.9,0.9,0.9}{
            \makecell[{{p{\linewidth}}}]{
                \texttt{\tiny{[GM$|$GM]}}
                \texttt{Expression: The grid that has a row with three X's in the first, third, and fifth columns of the second row, and a row with two X's in the first and third columns of the fourth row.} \\
            }
        }
    }
    & & \\ \\

    \theutterance \stepcounter{utterance}  
    & & \multicolumn{4}{p{0.6\linewidth}}{
        \cellcolor[rgb]{0.9,0.9,0.9}{
            \makecell[{{p{\linewidth}}}]{
                \texttt{\tiny{[GM$\rangle$P2]}}
                \texttt{You have attempted this same task before and did not succeed. These are the plans you wrote after those attempts.} \\
\texttt{Trial \#1: New plan: I failed to identify the correct grid because I didn't carefully examine the expression. The expression states that the grid has a row with three X's in the first, third, and fifth columns. However, I didn't check if the expression is referring to a row or a column. Next time, I will carefully read the expression and check if it's referring to a row or a column. I will also check each grid more thoroughly to see if it matches the expression.} \\
\\ 
\texttt{Trial \#2: New plan: I failed to identify the correct grid because I didn't carefully examine the expression. The expression states that the grid has a row with three X's in the first, third, and fifth columns. However, I didn't check if the expression is referring to a row or a column. Next time, I will carefully read the expression and check if it's referring to a row or a column. I will also check each grid more thoroughly to see if it matches the expression. Additionally, I will pay closer attention to the word "columns" in the expression, which implies that the X's are in the vertical direction, not the horizontal direction. This means I should be looking for rows with three X's in the first, third, and fifth positions, not columns.} \\
\\ 
\texttt{Trial \#3: New plan: I failed to identify the correct grid because I didn't carefully examine the expression. The expression states that the grid has a row with three X's in the first, third, and fifth columns, and a row with two □'s in the first and second columns. However, I didn't check if the expression is referring to a row or a column. Next time, I will carefully read the expression and check if it's referring to a row or a column. I will also check each grid more thoroughly to see if it matches the expression. Additionally, I will pay closer attention to the word "columns" in the expression, which implies that the X's are in the vertical direction, not the horizontal direction. This means I should be looking for rows with three X's in the first, third, and fifth positions, not columns. Furthermore, I will also check the rows with two □'s in the first and second columns, as specified in the expression.} \\
\\ 
\\ 
\texttt{You are given three grids, where each of them is 5 by 5 in size.} \\
\texttt{Grids have empty cells marked with "▢" and filled cells marked with "X".} \\
\texttt{You are also given a referring expression that describes one of the given grids.} \\
\texttt{Your task is to select the grid that matches the given referring expression.} \\
\texttt{Generate only the number (in text) of the grid that matches the given expression by selecting first, second, or third.} \\
\\ 
\texttt{First:} \\
\\ 
\texttt{□ X □ X □} \\
\texttt{X □ X □ X} \\
\texttt{□ X □ X □} \\
\texttt{X □ X □ X} \\
\texttt{□ X □ X □} \\
\\ 
\texttt{Second:} \\
\\ 
\texttt{X X □ X □} \\
\texttt{□ X X □ X} \\
\texttt{X □ □ X □} \\
\texttt{□ X □ □ X} \\
\texttt{X □ X □ X} \\
\\ 
\texttt{Third:} \\
\\ 
\texttt{X X □ □ X} \\
\texttt{□ X X □ X} \\
\texttt{X □ X X □} \\
\texttt{□ X □ X X} \\
\texttt{X □ □ X X} \\
\\ 
\texttt{Expression: The grid that has a row with three X's in the first, third, and fifth columns of the second row, and a row with two X's in the first and third columns of the fourth row.} \\
\texttt{Question: Which grid does the expression refer to?} \\
\texttt{Start with the tag "Answer: ", followed by your selection. Omit any other text.} \\
            }
        }
    }
    & \\ \\

    \theutterance \stepcounter{utterance}  
    & & & \multicolumn{4}{p{0.6\linewidth}}{
        \cellcolor[rgb]{0.9,0.9,0.9}{
            \makecell[{{p{\linewidth}}}]{
                \texttt{\tiny{[GM$\langle$P2]}}
                \texttt{Answer: second} \\
            }
        }
    }
     \\ \\

    \theutterance \stepcounter{utterance}  
    & & & \multicolumn{2}{p{0.3\linewidth}}{
        \cellcolor[rgb]{0.9,0.9,0.9}{
            \makecell[{{p{\linewidth}}}]{
                \texttt{\tiny{[GM$|$GM]}}
                \texttt{second} \\
            }
        }
    }
    & & \\ \\

    \theutterance \stepcounter{utterance}  
    & \multicolumn{4}{p{0.6\linewidth}}{
        \cellcolor[rgb]{0.95,0.95,0.95}{
            \makecell[{{p{\linewidth}}}]{
                \texttt{\tiny{[P2$\rightarrow$P1]}}
                \texttt{Post{-}episode Reflexion (not sent to GM):} \\
\texttt{New plan: I failed to distinguish the target grid from the distractors because my description was too general and didn't take into account the unique features of the target grid. The target grid has a row with three X's in the first, third, and fifth columns, but this description also applies to the first row of the first distractor grid. However, I didn't consider the arrangement of the □'s in the third row, which is unique to the target grid. Next time, I will focus on a more specific feature that is unique to the target grid, such as the arrangement of the □'s in the third row, and combine it with the row with three X's in the second row. For example, I could describe the target grid as "The grid that has a row with three X's in the first, third, and fifth columns of the second row, and a row with two □'s in the first and second columns of the third row."} \\
            }
        }
    }
    & & \\ \\

    \theutterance \stepcounter{utterance}  
    & & & \multicolumn{4}{p{0.6\linewidth}}{
        \cellcolor[rgb]{0.95,0.95,0.95}{
            \makecell[{{p{\linewidth}}}]{
                \texttt{\tiny{[P1$\rightarrow$P2]}}
                \texttt{Post{-}episode Reflexion (not sent to GM):} \\
\texttt{New plan: I failed to identify the correct grid because I didn't carefully examine the expression. The expression states that the grid has a row with three X's in the first, third, and fifth columns, and a row with two X's in the first and third columns of the fourth row. However, I didn't check if the expression is referring to a row or a column. Next time, I will carefully read the expression and check if it's referring to a row or a column. I will also check each grid more thoroughly to see if it matches the expression. Additionally, I will pay closer attention to the word "columns" in the expression, which implies that the X's are in the vertical direction, not the horizontal direction. This means I should be looking for rows with three X's in the first, third, and fifth positions, not columns. Furthermore, I will also check the rows with two X's in the first and third columns, as specified in the expression. I will also make sure to check the correct row for the second condition, which is the fourth row, not the second row as I previously did.} \\
            }
        }
    }
     \\ \\

\end{supertabular}
}

\end{document}
